
\documentclass[11pt, letterpaper]{article}


\usepackage[utf8]{inputenc}
\usepackage[spanish, es-tabla]{babel}
\usepackage[margin=2.3cm, top=2.8cm, bottom=2.8cm]{geometry}
\usepackage{helvet} 
\renewcommand{\familydefault}{\sfdefault}
\usepackage{microtype}
\usepackage{xcolor}
\usepackage{graphicx}
\usepackage{booktabs}
\usepackage{tabularx}
\usepackage{fancyhdr}
\usepackage{lastpage}
\usepackage{hyperref}
\usepackage{fontawesome5} % Para incluir el ícono de GitHub


\definecolor{usmblue}{RGB}{0, 75, 141}
\definecolor{usmred}{RGB}{200, 16, 46}
\definecolor{usmgray}{RGB}{100, 100, 100}
\definecolor{lightgray}{RGB}{245, 245, 245}


\pagestyle{fancy}
\fancyhf{}
\renewcommand{\headrulewidth}{0.4pt}
\renewcommand{\footrulewidth}{0.4pt}


\fancyhead[L]{\small \color{usmgray} Universidad Técnica Federico Santa María\\ \textbf{ICN-292 — Sistemas de Información para la Gestión}}
\fancyhead[R]{\small \color{usmgray} Departamento de Industrias\\ Segundo Semestre 2026}


\fancyfoot[L]{\small \color{usmgray} Laboratorio 5 · ICN-292 · Informe Oficial}
\fancyfoot[R]{\small \color{usmgray} Página \thepage\ de \pageref{LastPage}}

\hypersetup{
    colorlinks=true,
    linkcolor=usmblue,
    urlcolor=usmblue,
    pdftitle={Laboratorio 1 - ICN292}
}

\begin{document}


\thispagestyle{empty}


\begin{raggedright}
    {\bfseries\large UNIVERSIDAD TECNICA FEDERICO SANTA MARIA}\\[0.1cm]
    {\large Departamento de Industrias}\\[0.3cm]
    {\color{usmblue} \rule{\textwidth}{1.5pt}}\\[0.4cm]
    {\bfseries\color{usmblue}\huge INFORME DE LABORATORIO 5}\\[0.2cm]
    {\Large Entregable — SQL (2/2)
}\\[0.1cm]
    {\large\textbf{ICN-292 — Sistemas de Información para la Gestión}}\\[0.1cm]
    {\large Segundo Semestre 2026}
\end{raggedright}

\vspace{0.3cm}


\begin{center}
\colorbox{lightgray}{
\begin{minipage}{0.95\textwidth}
\vspace{0.2cm}
\begin{tabular}{ll}
    \textbf{Estudiante:} & Martin Ariel Aspee Arroyo \\
   
    \textbf{Fecha de Entrega:} & Sábado 14 de Noviembre de 2026, 11:00 hrs \\
    \textbf{Campus / Sede:} & Vitacura \\
    \textbf{Profesores:} & José Miguel González / Paul José Luis Sáez Tamayo
\end{tabular}
\vspace{0.2cm}
\end{minipage}
}
\end{center}


\vspace{0.3cm}
\begin{center}
    \href{https://github.com/Diaspuu/ICN292-Lab5-Aspee-Martin}{\Large \faGithub} \\
    \vspace{0.1cm}
    \url{https://github.com/Diaspuu/ICN292-Lab5-Aspee-Martin}
\end{center}
\vspace{0.3cm}
\section{Diagrama Entidad-Relación}
\begin{figure}[h]
    \centering
    \includegraphics[width=0.5\linewidth]{diagrama e-r.png}
    \caption{Modelo Entidad-Relación de la base de datos andesbike.}
    \label{fig:placeholder}
\end{figure}
\newpage
\section{consultas}
\subsection{consulta 1}
\begin{figure}[h]
    \centering
    \includegraphics[width=0.6\linewidth]{consulta 1.png}
    \caption{Vendedores y su cantidad de ventas, incluyendo aquellos sin registros.}
    \label{fig:placeholder}
\end{figure}

\subsection{consulta 2}
\begin{figure}[h]
    \centering
    \includegraphics[width=0.6\linewidth]{consulta 2.png}
    \caption{Tiendas con sus ventas y unidades totales, incluyendo tiendas sin vendedores.}
    \label{fig:placeholder}
\end{figure}
\newpage
\subsection{consulta 3}
\begin{figure}[h]
    \centering
    \includegraphics[width=0.6\linewidth]{consulta 3.png}
    \caption{Listado de tiendas que no tienen ningún vendedor asignado.}
    \label{fig:placeholder}
\end{figure}

\subsection{consulta 4}
\begin{figure}[h]
    \centering
    \includegraphics[width=0.6\linewidth]{consulta 4.png}
    \caption{Modelos de bicicletas sin ventas registradas en el mes de junio de 2025.}
    \label{fig:placeholder}
\end{figure}

\noindent\textbf{Explicación:} La condición del mes va en la cláusula ON del LEFT JOIN para intentar emparejar solo las ventas de junio. Si se hubiera puesto en el WHERE, se habría filtrado la tabla final completa, eliminando del resultado a los modelos sin ventas que se pedía mostrar.
\newpage
\subsection{consulta 5}
\begin{figure}[h]
    \centering
    \includegraphics[width=0.6\linewidth]{consulta 5.png}
    \caption{Vendedores que no han registrado ventas, indicando su tienda asociada.}
    \label{fig:placeholder}
\end{figure}

\subsection{consulta 6}
\begin{figure}[h]
    \centering
    \includegraphics[width=0.6\linewidth]{consulta 6.png}
    \caption{Pares de vendedores que comparten local, ordenados por antigüedad de ingreso.}
    \label{fig:placeholder}
\end{figure}
\newpage
\subsection{consulta 7}
\begin{figure}[h]
    \centering
    \includegraphics[width=0.6\linewidth]{consulta 7.png}
    \caption{Monto total vendido en categorías Montaña y Ruta para tiendas sobre \$2.000.000.}
    \label{fig:placeholder}
\end{figure}
\subsection{cosulta 8}
\begin{figure}[h]
    \centering
    \includegraphics[width=0.5\linewidth]{consulta 8.png}
    \caption{Bicicletas cuyo precio supera el promedio general del catálogo.}
    \label{fig:placeholder}
\end{figure}

\noindent\textbf{Explicación:} El valor del precio promedio del catálogo es 637490.00 y se obtuvo ejecutando una subconsulta de agregación (\texttt{SELECT AVG(precio) FROM bicicletas;}) directamente al motor de la base de datos.
\newpage
\subsection{consulta 9}
\begin{figure}[h]
    \centering
    \includegraphics[width=0.5\linewidth]{consulta 9.png}
    \caption{Vendedores cuyo promedio de unidades supera el promedio de ventas de toda la cadena.}
    \label{fig:placeholder}
\end{figure}

\subsection{consulta 10}
\begin{figure}[h]
    \centering
    \includegraphics[width=0.4\linewidth]{consulta 10.png}
    \caption{Monto total recaudado por cada vendedor, incluyendo aquellos sin ventas.}
    \label{fig:placeholder}
\end{figure}

\noindent\textbf{Explicación:} Algunos vendedores no muestran un cero en sus montos, sino que aparecen en blanco (\texttt{NULL}) porque la función matemática \texttt{SUM}, al intentar sumar sobre filas que el \texttt{LEFT JOIN} dejó vacías, devuelve un vacío matemático. Para que mostraran cero, tendría que agregar la función \texttt{COALESCE} envolviendo el cálculo, de la siguiente forma: \texttt{COALESCE(SUM(ve.unidades * b.precio), 0)}.

\end{document}